% Extraído de Textuais/03-Metodologia.tex (Seção de indicadores de desempenho no nível de método).
\begin{table}[htbp]
    \centering
    \footnotesize
    \renewcommand{\arraystretch}{1.2}
    \captionsetup{font+={color=ForestGreen}}%
    \caption{Indicadores de desempenho no nível de método}
    \label{tab:met-indicadores}
    \color{ForestGreen}%
    \begin{tabularx}{\linewidth}{@{}>{\raggedright\arraybackslash}p{2.6cm}>{\raggedright\arraybackslash}p{1.9cm}>{\raggedright\arraybackslash}X>{\raggedright\arraybackslash}p{3.3cm}@{}}
        \toprule
        Indicador & Papel & O que mede & Valores possíveis e valor de um classificador aleatório \\
        \midrule
        PR-AUC (Equação~\ref{eq:ap}) & Principal & Área sob a curva de precisão e \textit{recall}, que resume, em todos os limiares, quanto das previsões de mudança está correto à medida que mais mudanças são encontradas & 0 a 1; aleatório igual à taxa de positivos ($\approx 0{,}005$) \\
        ROC-AUC & Principal & Probabilidade de um método que muda receber do modelo uma probabilidade maior do que um método que não muda & 0 a 1; aleatório igual a 0,5 \\
        \midrule
        Precisão (Equação~\ref{eq:precisao}) & Complementar & Fração das previsões de mudança que estavam corretas & 0 a 1 \\
        \textit{Recall} (Equação~\ref{eq:recall}) & Complementar & Fração das mudanças reais que o modelo previu & 0 a 1 \\
        F1 da classe positiva (Equação~\ref{eq:f1score}) & Complementar & Média harmônica da precisão e do \textit{recall} & 0 a 1 \\
        Acurácia (Equação~\ref{eq:acuracia}) & Complementar & Fração de todas as previsões que estavam corretas & 0 a 1; prever sempre ``não muda'' resulta em $\approx 0{,}995$ \\
        \bottomrule
    \end{tabularx}
    \fonte{Elaborado pelo autor (2026).}
\end{table}%
